% GENERATED FILE. Edit canonical lessons, then run npm run generate:books.
% canonical-source-hash: fnv1a64:f8ac3e85a51abbab
% canonical-lessons: HI-C415-mainejar, HI-C415-rajistar, HI-C415-buking, HI-C415-thikana, HI-C415-magana

\chapter{At the Front Desk}
\label{ch:hi-at-the-front-desk}
\hlchaptermodality{415}

\begin{chapteropening}
\textbf{By the end of this chapter:} \emph{I can speak to the manager, sign the register and send for tea.}

Complete the last lesson of chapter 415: I can speak to the manager, sign the register and send for tea.
\end{chapteropening}

\section[\texorpdfstring{\hi{मैनेजर} (mainejar)}{मैनेजर (mainejar)}]{\hi{मैनेजर} (mainejar) — a manager}

\label{lesson:HI-C415-mainejar}



\begin{quote}
Before the new one: say the Hindi for comfortable, then the Hindi for to wake (someone).
\end{quote}

\subsection*{You'll want to know: \hi{मैनेजर}}
\textbf{\hi{मैनेजर}} — \emph{mainejar} — \textquotedblleft{}a manager\textquotedblright{}.

\begin{quote}
\hi{होटल} \hi{का} \hi{मैनेजर} \hi{बहुत} \hi{अच्छा} \hi{आदमी} \hi{था।}
\end{quote}

\begin{quote}
\emph{hoṭal kā mainejar bahut acchā ādmī thā.}
\end{quote}

\begin{quote}
The hotel manager was a very nice man.
\end{quote}

Masculine or feminine, by the person meant.

\begin{cousinweb}
\textbf{\hi{मैनेजर}} (\emph{mainejar}) is borrowed from English \emph{manager}.
\end{cousinweb}

\begin{grammarlens}[title={What you've built}]
One more word for getting around.
\end{grammarlens}

\subsection*{Your turn}
\begin{itemize}
\raggedright
  \item \emph{Say it:} \emph{mainejar}
  \item \emph{Say it:} \emph{mainejar}, once more
  \item \emph{Say it:} \emph{hoṭal kā mainejar bahut acchā ādmī thā.} — The hotel manager was a very nice man.
  \item \emph{Recall:} say the Hindi for comfortable, then the Hindi for to wake (someone), then say \emph{mainejar} again
\end{itemize}

\begin{tcolorbox}[breakable,colback=teal!4,colframe=teal!35!black,title={Before you move on}]
What does \hi{मैनेजर} mean? (\textquotedblleft{}A manager\textquotedblright{}.) Say it once more.
\end{tcolorbox}
\section[\texorpdfstring{\hi{रजिस्टर} (rajisṭar)}{रजिस्टर (rajisṭar)}]{\hi{रजिस्टर} (rajisṭar) — a register (book)}

\label{lesson:HI-C415-rajistar}



\begin{quote}
Before the new one: say the Hindi for to wake (someone), then the Hindi for a manager.
\end{quote}

\subsection*{You'll want to know: \hi{रजिस्टर}}
\textbf{\hi{रजिस्टर}} — \emph{rajisṭar} — \textquotedblleft{}a register (book)\textquotedblright{}.

\begin{quote}
\hi{मैनेजर} \hi{ने} \hi{रजिस्टर} \hi{में} \hi{मेरा} \hi{नाम} \hi{लिखा।}
\end{quote}

\begin{quote}
\emph{mainejar ne rajisṭar meṁ merā nām likhā.}
\end{quote}

\begin{quote}
The manager wrote my name in the register.
\end{quote}

Masculine.

\begin{cousinweb}
\textbf{\hi{रजिस्टर}} (\emph{rajisṭar}) is borrowed from English \emph{register}, the book at a hotel desk where every guest's name goes.
\end{cousinweb}

\begin{grammarlens}[title={What you've built}]
One more word for getting around.
\end{grammarlens}

\subsection*{Your turn}
\begin{itemize}
\raggedright
  \item \emph{Say it:} \emph{rajisṭar}
  \item \emph{Say it:} \emph{rajisṭar}, once more
  \item \emph{Say it:} \emph{mainejar ne rajisṭar meṁ merā nām likhā.} — The manager wrote my name in the register.
  \item \emph{Recall:} say the Hindi for to wake (someone), then the Hindi for a manager, then say \emph{rajisṭar} again
\end{itemize}

\begin{tcolorbox}[breakable,colback=teal!4,colframe=teal!35!black,title={Before you move on}]
What does \hi{रजिस्टर} mean? (\textquotedblleft{}A register (book)\textquotedblright{}.) Say it once more.
\end{tcolorbox}
\section[\texorpdfstring{\hi{बुकिंग} (buking)}{बुकिंग (buking)}]{\hi{बुकिंग} (buking) — a booking}

\label{lesson:HI-C415-buking}



\begin{quote}
Before the new one: say the Hindi for a manager, then the Hindi for a register.
\end{quote}

\subsection*{You'll want to know: \hi{बुकिंग}}
\textbf{\hi{बुकिंग}} — \emph{buking} — \textquotedblleft{}a booking\textquotedblright{}.

\begin{quote}
\hi{मेरी} \hi{बुकिंग} \hi{तीन} \hi{रात} \hi{की} \hi{है।}
\end{quote}

\begin{quote}
\emph{merī buking tīn rāt kī hai.}
\end{quote}

\begin{quote}
My booking is for three nights.
\end{quote}

Feminine.

\begin{cousinweb}
\textbf{\hi{बुकिंग}} (\emph{buking}) is borrowed from English \emph{booking}; it is feminine in Hindi.
\end{cousinweb}

\begin{grammarlens}[title={What you've built}]
One more word for getting around.
\end{grammarlens}

\subsection*{Your turn}
\begin{itemize}
\raggedright
  \item \emph{Say it:} \emph{buking}
  \item \emph{Say it:} \emph{buking}, once more
  \item \emph{Say it:} \emph{merī buking tīn rāt kī hai.} — My booking is for three nights.
  \item \emph{Recall:} say the Hindi for a manager, then the Hindi for a register, then say \emph{buking} again
\end{itemize}

\begin{tcolorbox}[breakable,colback=teal!4,colframe=teal!35!black,title={Before you move on}]
What does \hi{बुकिंग} mean? (\textquotedblleft{}A booking\textquotedblright{}.) Say it once more.
\end{tcolorbox}
\section[\texorpdfstring{\hi{ठिकाना} (ṭhikānā)}{ठिकाना (ṭhikānā)}]{\hi{ठिकाना} (ṭhikānā) — a place to stay; whereabouts}

\label{lesson:HI-C415-thikana}



\begin{quote}
Before the new one: say the Hindi for a register, then the Hindi for a booking.
\end{quote}

\subsection*{You'll want to know: \hi{ठिकाना}}
\textbf{\hi{ठिकाना}} — \emph{ṭhikānā} — \textquotedblleft{}a place to stay; whereabouts\textquotedblright{}.

\begin{quote}
\hi{रात} \hi{के} \hi{लिए} \hi{हमें} \hi{कोई} \hi{ठिकाना} \hi{नहीं} \hi{मिला।}
\end{quote}

\begin{quote}
\emph{rāt ke liye hamẽ koī ṭhikānā nahīṃ milā.}
\end{quote}

\begin{quote}
We found nowhere to stay for the night.
\end{quote}

Masculine.

\begin{cousinweb}
\textbf{\hi{ठिकाना}} (\emph{ṭhikānā}) goes back to \emph{*ṭhīkka}, a form CDIAL reconstructs (CDIAL 5503), the same entry as \hi{ठीक}; it lists Hindi \emph{ṭhikānā}, \textquotedblleft{}a station\textquotedblright{}. A \hi{ठिकाना} is a fixed place where you can be found.
\end{cousinweb}

\begin{grammarlens}[title={What you've built}]
One more word for getting around.
\end{grammarlens}

\subsection*{Your turn}
\begin{itemize}
\raggedright
  \item \emph{Say it:} \emph{ṭhikānā}
  \item \emph{Say it:} \emph{ṭhikānā}, once more
  \item \emph{Say it:} \emph{rāt ke liye hamẽ koī ṭhikānā nahīṃ milā.} — We found nowhere to stay for the night.
  \item \emph{Recall:} say the Hindi for a register, then the Hindi for a booking, then say \emph{ṭhikānā} again
\end{itemize}

\begin{tcolorbox}[breakable,colback=teal!4,colframe=teal!35!black,title={Before you move on}]
What does \hi{ठिकाना} mean? (\textquotedblleft{}A place to stay; whereabouts\textquotedblright{}.) Say it once more.
\end{tcolorbox}
\section[\texorpdfstring{\hi{मँगाना} (mãgānā)}{मँगाना (mãgānā)}]{\hi{मँगाना} (mãgānā) — to send for; to order (to be brought)}

\label{lesson:HI-C415-magana}



\begin{quote}
Before the new one: say the Hindi for a booking, then the Hindi for a place to stay.
\end{quote}

\subsection*{You'll want to know: \hi{मँगाना}}
\textbf{\hi{मँगाना}} — \emph{mãgānā} — \textquotedblleft{}to send for; to order (to be brought)\textquotedblright{}.

\begin{quote}
\hi{हमने} \hi{कमरे} \hi{में} \hi{चाय} \hi{मँगाई।}
\end{quote}

\begin{quote}
\emph{hamne kamre meṁ cāy mãgāī.}
\end{quote}

\begin{quote}
We ordered tea to our room.
\end{quote}

Past, with \hi{ने} (he or she …): \hi{उसने} \hi{खाना} \hi{मँगाया}, \emph{usne khānā mãgāyā}.

\begin{cousinweb}
\textbf{\hi{मँगाना}} (\emph{mãgānā}) is the causative partner of \hi{माँगना}, \textquotedblleft{}to ask for\textquotedblright{}, which goes back to Sanskrit \emph{mārgati}, \textquotedblleft{}seeks, requests\textquotedblright{} (CDIAL 10074): you ask for a thing to be brought to you.
\end{cousinweb}

\begin{grammarlens}[title={What you've built}]
One more word for getting around.
\end{grammarlens}

\subsection*{Your turn}
\begin{itemize}
\raggedright
  \item \emph{Say it:} \emph{mãgānā}
  \item \emph{Say it:} \emph{mãgānā}, once more
  \item \emph{Say it:} \emph{hamne kamre meṁ cāy mãgāī.} — We ordered tea to our room.
  \item \emph{Recall:} say the Hindi for a booking, then the Hindi for a place to stay, then say \emph{mãgānā} again
\end{itemize}

\begin{tcolorbox}[breakable,colback=teal!4,colframe=teal!35!black,title={Before you move on}]
What does \hi{मँगाना} mean? (\textquotedblleft{}To send for; to order (to be brought)\textquotedblright{}.) Say it once more.
\end{tcolorbox}
